\section{Conclusion}

We have identified a specific, correctable mismatch in MoE mixed-precision quantization: sensitivity-based allocators treat every expert's reconstruction loss equally, ignoring that routing frequencies vary by orders of magnitude. Route-Weighted Sensitivity corrects this by scaling each expert's sensitivity by its routing weight before the ILP---a modification to the objective that is principled (approximating the expected inference loss), conservative (bounded $2\times$ scaling), and effective (consistent PPL improvements on all three models at aggressive budgets). The method requires no tuning and no additional calibration cost. The Route Heterogeneity Index ($\rhi$) provides a cheap preliminary indicator of whether routing-aware allocation is worth trying.

Our cross-layer ablation shows that naive propagation of routing information fails catastrophically, supporting the design choice of local, per-layer route weighting. The broader lesson is that \emph{the allocator objective should reflect how often a weight is used at inference}; future MoE compression work could extend this principle to latency profiles, outlier smoothing, and activation-aware sensitivity.

\paragraph{Limitations.}
Our claims are deliberately scoped: \method\ is evaluated with a binary strategy pool and RTN-based per-block quantization; generalization to larger strategy pools and GPTQ/AWQ-based methods requires testing. The $(1 + w)$ weighting is a conservative heuristic; systematic comparison with raw-$\piprob$, gate-mass, and conditional-loss weighting across all models would further validate the design. RHI is presented as a preliminary diagnostic with only three data points; additional models are needed to confirm its utility. Calibration-subset sensitivity across 3 seeds confirms consistent improvements (Table~\ref{tab:seed}), though more seeds and diverse calibration sources would strengthen robustness claims.
